\scanpage{44}
$G$, groupe semi-simple adjoint sur base $S$, de type constant.

(i) À tout couple de Killing $(T,B)$ dans $G$, on associe un hom. (en fait un mono)
\[
G_{\mathrm m}\overset{i}{\lhook\joinrel\longrightarrow}G
\]
par descente à partir du cas où le système de racines sur $T$ est déployé, donc où $T\xrightarrow{\sim}G_{\mathrm m}^{\Sigma}$ ($\Sigma$ ens des racines simples), \deleted{\unsure{on prend|quand}} d'où $\mathrm{Hom}(G_{\mathrm m},G_{\mathrm m})\simeq\mathbb Z^{\Sigma}$ : on prend l'élément dont toutes les coordonnées sont $1$.

Le centralisateur d'un tel hom (i.e. d'un tel sous-tore) \deleted{est} $T$, donc le $\unsure{=|\sim}$ que pour le couple de Killing.

On appelle \uline{hom principal} $i$ de $G_{\mathrm m}$ dans $G$ un hom tel qu'il existe un couple de Killing auquel il est associé. Ce couple est alors unique [$i$ est régulier], et par suite la condition envisagée sur $i$ est loc\textsuperscript{le} pour fpqc. Il y a corr. biunivoque entre hom. principaux $G_{\mathrm m}\unsure{-|\to}G$ et couples de Killing.

(i bis) \deleted{Le sous-tore normalisateur de $i(G_{\mathrm m})$, $i$ principal, est} Si $i$ est principal, $i^{-1}=[\lambda\mapsto i(\lambda^{-1})]$ aussi, il est associé au couple de Killing opposé. Le normalisateur de $i(G_{\mathrm m})$ est le ss-groupe\inserted{$N_{\unsure{1|0}}(T)$} de $N(T)$ image inverse du sous-groupe $\simeq(\mathbb Z/2\mathbb Z)_S\subset W$ engendré par l'élément $w_{(T,B)}$ de $W$ qui envoie $B$ sur le Borel opposé.
\margin{$r\geq 1$ !}

Sous-tore de rang $1$ principal : qui loc. fpqc est \deleted{\unsure{de}}l'image d'un hom principal $\Longleftrightarrow$ qui sur un rev. étale \deleted{\unsure{galoisien}} principal de \deleted{groupe} $\mathbb Z/2\mathbb Z$ est comme \unsure{ça}.
% archive sheet number: 43

